%%% --- CORPO DO TRABALHO --- %%%
\chapter{Metodologia}
\label{c2_metodologias_tg}
%\thispagestyle{empty}

Este trabalho busca investigar os efeitos dos riscos climáticos sobre a estabilidade financeira, com foco específico na relação entre a ocorrência de desastres ambientais e o comportamento da inadimplência no sistema de crédito bancário. Para tanto, adota-se uma abordagem quantitativa baseada em séries temporais, estruturada em três etapas metodológicas  descritas nas seções seguintes.

\section{Causalidade de Granger}
A primeira etapa metodológica consiste na aplicação de testes de causalidade de Granger, com o objetivo de identificar relações de precedência temporal entre os indicadores de desastres ambientais e as taxas de inadimplência das carteiras de crédito bancário. O conceito de causalidade de Granger estabelece que uma variável $X$ causa, no sentido de Granger, uma variável $Y$ se a inclusão dos valores passados de $X$ melhora significativamente a previsão dos valores futuros de $Y$, em comparação a um modelo que utiliza apenas os valores defasados da própria variável $Y$ \citep{granger1969}. Formalmente, dado um modelo autorregressivo vetorial (VAR), testa-se a hipótese nula de que os coeficientes associados às defasagens de $X$ são conjuntamente nulos na equação de $Y$, ou seja, de que $X$ não contém informação preditiva adicional sobre $Y$ além daquela já capturada por seu próprio histórico.
Nesse contexto, essa abordagem permite investigar o poder explicativo de eventos climáticos passados sobre a inadimplência, bem como testar a existência de precedência temporal entre a ocorrência de desastres ambientais e a dinâmica do sistema de crédito nacional.


\subsection{Teste de Granger}
A ideia por trás do teste consiste na comparação entre dois modelos de regressão linear. O \textbf{modelo restrito}, assim denominado pois restringe a zero a constribuição  de $X$, que corresponde a um modelo autorregressivo de $Y$ de ordem $p$, na qual apenas o histórico da própria série é empregado como preditor, é definido como

\begin{equation}
    Y_t \;=\; \alpha_0 \;+\; \sum_{j=1}^{p} \alpha_j\, Y_{t-j} \;+\; u_t,
    \label{eq:restrito}
\end{equation}

\noindent em que:
\begin{itemize}
    \item $Y_t$ é o valor da variável dependente no período $t$;
    \item $\alpha_0$ é o intercepto do modelo;
    \item $\alpha_j$ são os coeficientes associados às $p$ defasagens de $Y$;
    \item $p$ é o número de defasagens \textit{(lags)} incluídas no modelo;
    \item $u_t$ é um ruído branco.
\end{itemize}

A estimação desse modelo por Mínimos Quadrados Ordinários (MQO) fornece a \textbf{Soma dos Quadrados dos Resíduos Restrita} ($\mathrm{SQR}_r$), que quantifica o erro de previsão de $Y$ na ausência da variável $X$.

Caso $X$ contenha informação adicional relevante para a previsão de $Y$, além daquela já capturada por seus próprios valores passados, a inclusão das defasagens de $X$ deve reduzir o erro de previsão. Essa ideia é formalizada no \textbf{modelo irrestrito}, que incorpora simultaneamente as defasagens de $Y$ e de $X$ através de

\begin{equation}
    Y_t \;=\; \beta_0 \;+\; \sum_{j=1}^{p} \beta_j\, Y_{t-j} \;+\; \sum_{j=1}^{p} \delta_j\, X_{t-j} \;+\; \varepsilon_t,
    \label{eq:irrestrito}
\end{equation}

\noindent em que:
\begin{itemize}
    \item $Y_t$ é o valor da variável dependente no período $t$;
    \item $\beta_0$ é o intercepto do modelo;
    \item $\beta_j$ são os coeficientes associados às $p$ defasagens de $Y$;
    \item $\delta_j$ são os coeficientes associados às $p$ defasagens de $X$;
    \item $p$ é o número de defasagens \textit{(lags)} incluídas no modelo;
    \item $\varepsilon_t$ é o termo de erro, assumido como ruído branco.
\end{itemize}

A estimação por Mínimos Quadrados Ordinários (MQO) fornece a \textbf{Soma dos Quadrados dos Resíduos Irrestrita} ($\mathrm{SQR}_u$). 

A ideia central do teste consiste em verificar se os coeficientes das defasagens de $X$ no modelo irrestrito são conjuntamente significativos. Em outras palavras, devemos testar as hipóteses:

\begin{equation}
\left\{
\begin{array}{ll}
H_0: & \delta_1 = \delta_2 = \cdots = \delta_p = 0 \\[10pt]
H_1: & \delta_j \neq 0, \quad \textnormal{para pelo menos um } j \in \{1, \ldots, p\},
\end{array}
\right.
\label{eq:hipoteses}
\end{equation}

\noindent em que $H_0$ estabelece que os coeficientes associados às $p$ defasagens de $X$ são conjuntamente nulos no modelo irrestrito, o que implica que $X$ não contribui significativamente para melhorar a capacidade preditiva de $Y$. Enquanto, $H_1$ sugere que pelo menos um dos coeficientes $\delta_j$ é estatisticamente diferente de zero, indicando que os valores passados de $X$ possuem informação relevante para a previsão de $Y$.


As hipóteses são testadas por meio de uma estatística $F$, obtida pela comparação entre $\mathrm{SQR}_r$ e $\mathrm{SQR}_u$. Se  $X$ de fato causa, no sentido de Granger, $Y$, a inclusão de suas defasagens deve reduzir o erro de previsão de forma significativa, isto é, $\mathrm{SQR}_u < \mathrm{SQR}_r$, implicando valores elevados para a estatística $F$. Logo, a estatística de teste é dada por

\begin{equation}
    F \;=\; \frac{(\mathrm{SQR}_r - \mathrm{SQR}_u)\,/\,p}{\mathrm{SQR}_u\,/\,(n - k)} \;\sim\; F_{(p,\; n-k)},
    \label{eq:estatistica_F}
\end{equation}

\noindent em que:
\begin{itemize}
    \item $\mathrm{SQR}_r$ representa a soma dos quadrados dos resíduos do \textbf{modelo restrito}, obtida pela regressão de $Y_t$ apenas sobre suas próprias defasagens;
    
    \item $\mathrm{SQR}_u$ representa a soma dos quadrados dos resíduos do \textbf{modelo irrestrito}, obtida pela regressão de $Y_t$ sobre suas defasagens e sobre as defasagens de $X_t$;
    
    \item $\mathrm{SQR}_r - \mathrm{SQR}_u$ corresponde à redução no erro de ajuste decorrente da inclusão das defasagens de $X$; quanto maior essa diferença, maior a evidência de que $X$ contribui para a previsão de $Y$;
    
    \item $p$ denota o número de restrições impostas, equivalente ao número de defasagens de $X$ incluídas no modelo irrestrito, e determina os graus de liberdade do numerador;
    
    \item $n$ denota o número total de observações da amostra;
    
    \item $k$ denota o número total de parâmetros estimados no modelo irrestrito (intercepto mais $p$ coeficientes associados a $Y$ e $p$ coeficientes associados a $X$, de modo que $k = 2p + 1$);
    
    \item $n - k$ corresponde aos graus de liberdade associados ao denominador da estatística $F$.
\end{itemize}

\noindent Sob $H_0$, a estatística F segue assintoticamente uma distribuição $F$ com $p$ e $n - k$ graus de liberdade.

Se o valor-$p$ associado à estatística $F$ for inferior ao nível de significância $\alpha$, rejeita-se $H_0$, concluindo-se que há evidências de que $X$ Granger-causa $Y$. Em outras palavras, os valores passados de $X$ contêm informação estatisticamente relevante para a previsão de $Y$, além daquela já incorporada em seu próprio passado. Caso contrário, não se rejeita $H_0$, não havendo evidências suficientes de causalidade de Granger na direção $X \to Y$.


Ressalta-se que a causalidade de Granger é, em geral, assimétrica, isto é, a rejeição de $H_0$ na direção $X \to Y$ não implica, necessariamente, que $Y$ também cause, no sentido de Granger, $X$.

A validade do teste de Granger está condicionada a um conjunto de suposições que devem ser verificadas previamente. A primeira é a de \textbf{linearidade}. O teste assume que as relações entre $X$ e $Y$ são lineares, conforme especificado nos modelos restrito e irrestrito. Relações não lineares entre as variáveis não são captadas pelo procedimento padrão.

Além disso, requer-se \textbf{estacionaridade} das séries, isto é, as séries temporais envolvidas devem apresentar média, variância e autocovariâncias constantes ao longo do tempo. Na presença de raiz unitária (comportamento não-estacionário), a aplicação direta do teste pode levar a inferências inválidas \citep{phillips1988}. Nesses casos, recomenda-se testar a estacionaridade,  por exemplo, via teste de Dickey Fuller Aumentado \citep{dickey1979} ou via teste de \cite{phillips1988}. Caso as séries sejam identificadas como não estacionárias, emprega-se procedimentos alternativos, como o de \citet{toda1995}, que viabiliza a aplicação do teste de Granger mesmo na presença de raiz unitária e/ou cointegração, isto é, mesmo quando as séries compartilham uma tendência estocástica comum de longo prazo. O método consiste em estimar um modelo $\mathrm{VAR}(p + d_{\max})$, onde $d_{\max}$ é a ordem máxima de integração das séries, e aplicar o teste de Wald apenas sobre os primeiros $p$ coeficientes, garantindo a validade assintótica da inferência sem a necessidade de pré-testes de cointegração.


Ademais, vale ressaltar que o teste é sensível à escolha do número de defasagens $p$. Uma especificação inadequada pode comprometer os resultados, sendo comum selecionar $p$ com base em critérios de informação, como o Critério de Informação de Akaike (AIC) ou o Critério Bayesiano de Schwarz (BIC).


Por fim, ressalta-se que a causalidade de Granger é um conceito de precedência temporal e previsibilidade, e não de causalidade estrutural. A rejeição de $H_0$ indica que $X$ antecede temporalmente $Y$ e melhora sua capacidade de previsão, mas não implica, por si só, a existência de uma relação causal estrutural entre as variáveis.



\section{Modelos de Defasagens Distribuídas}

Dados de séries temporais frequentemente exibem efeitos defasados, em que valores passados
de uma variável influenciam o comportamento presente de outra, e modelos dinâmicos capturam
essas dependências ao incorporar explicitamente valores defasados das variáveis explicativas
\citep{almon1965}. No contexto deste trabalho, essa característica é especialmente relevante, pois
eventos climáticos raramente impactam a inadimplência de forma imediata. Uma seca severa,
por exemplo, tende a reduzir a produção agrícola ao longo dos meses subsequentes, de modo
que o produtor rural pode entrar em inadimplência apenas no vencimento de suas obrigações
de crédito, o que pode ocorrer com uma defasagem de três a seis meses em relação ao evento.
De forma semelhante, uma enchente que destrói domicílios e interrompe atividades econômicas
pode afetar a capacidade de pagamento das famílias em prazos mais curtos, gerando aumento
da inadimplência em um a dois meses.

Diante dessa dinâmica, estimam-se modelos de defasagens distribuídas (\textit{Distributed Lag
Models} - DLM), que permitem capturar como os efeitos dos choques climáticos se propagam ao
longo do tempo sobre a inadimplência, reconhecendo que tais impactos podem não ser imediatos,
mas manifestar-se de forma gradual e persistente nos períodos subsequentes ao evento.

\subsection{Definição e Formulação do Modelo}

Um modelo de defasagens distribuídas relaciona uma variável dependente $Y_t$ a valores
correntes e passados de uma variável explicativa $X_t$, permitindo que o efeito de $X$ sobre $Y$
se distribua ao longo de múltiplos períodos \citep{almon1965}. Na forma geral, o modelo é
especificado como

\begin{equation}
    Y_t \;=\; \alpha \;+\; \sum_{k=0}^{K} \beta_k\, X_{t-k} \;+\; \varepsilon_t,
    \label{eq:dlm_geral}
\end{equation}

\noindent em que:

\begin{itemize}
    \item $Y_t$ é o valor da variável dependente no período $t$, no contexto do trabalho, a taxa de inadimplência;
    \item $\alpha$ é o intercepto do modelo;
    \item $X_{t-k}$ é o valor da variável explicativa com $k$ períodos de defasagem. Aqui, o indicador de desastres ambientais ou climáticos;
    \item $\beta_k$ é o coeficiente associado à defasagem $k$, que mede o efeito de $X$ sobre $Y$ após $k$ períodos;
    \item $K$ é o número máximo de defasagens incluídas no modelo;
    \item $\varepsilon_t$ é o termo de erro aleatório, assumido com distribuição gaussiana
          de média zero e variância constante, que captura a variação de $Y_t$ não
          explicada pelas defasagens de $X$ \citep{liang2024}.
\end{itemize}

Cada coeficiente $\beta_k$ representa o \textbf{efeito parcial} de $X_{t-k}$ sobre $Y_t$,
isto é, a variação esperada em $Y_t$ decorrente de uma variação unitária em $X$ ocorrida $k$
períodos antes, mantidos constantes os demais termos. O conjunto dos coeficientes
$\{\beta_0, \beta_1, \ldots, \beta_K\}$ forma o chamado \textbf{perfil de resposta defasada}, que descreve a trajetória temporal do impacto de $X$ sobre $Y$
\citep{wooldridge2010}.

Duas quantidades são de especial interesse. O \textbf{efeito de impacto}, que corresponde a
$\beta_0$ e mede a resposta imediata de $Y_t$ a uma variação em $X_t$, no mesmo período.
A \textbf{associação acumulada no horizonte $K$} é dada pela soma de todos os coeficientes:

\begin{equation}
    B_K \;=\; \sum_{k=0}^{K} \beta_k,
    \label{eq:efeito_lp}
\end{equation}

\noindent No DLM em níveis, essa soma descreve a mudança prevista de $Y$ após a transmissão de uma mudança permanente unitária em $X$, sob a especificação adotada. Na aplicação, porém, as variáveis são $y_{it}=\Delta\log(I_{it})$ e $x_{it}=\Delta\log(1+D_{it})$, sendo $I$ a taxa de inadimplência e $D$ a contagem administrativa de registros. A soma $B_K$ descreve a associação acumulada em $\log(I)$, até $K$, a um pulso unitário em $x$; não representa automaticamente o impacto total de um desastre físico nem pontos percentuais da taxa.

Uma mudança persistente $D_0\to D_1$ gera um pulso $\delta=\log(1+D_1)-\log(1+D_0)$ em $x$. Sua resposta acumulada até $K$ é $\delta B_K$ em unidades logarítmicas, equivalente a $100[\exp(\delta B_K)-1]$ por cento de variação relativa de $I$. Uma mudança de $D$ por apenas um mês produz dois pulsos, $\delta$ e $-\delta$, cuja resposta requer convolução com o perfil de coeficientes. Essas interpretações são associações sob cenários definidos, sem identificação causal. Somatórios com valores distintos de $K$ se referem a horizontes distintos.

\subsection{Modelo com Efeitos Fixos}

Quando os dados possuem estrutura em painel, isto é, quando várias unidades $i$
são acompanhadas ao longo do tempo $t$, surge uma dificuldade recorrente:
características próprias de cada unidade, que permanecem relativamente constantes
no tempo, podem influenciar simultaneamente a variável dependente e as variáveis
explicativas. Nessa situação, caso tais fatores não sejam considerados, parte do
efeito atribuído às variáveis de interesse poderá, na realidade, refletir diferenças
estruturais entre as unidades.

No contexto deste trabalho, por exemplo, um estado historicamente mais vulnerável
a eventos climáticos extremos pode registrar, ao mesmo tempo, maior incidência de
desastres naturais e uma atividade econômica mais sensível a restrições de crédito.
Se essa heterogeneidade não observada for ignorada, os coeficientes $\beta_k$
estimados tenderão a captar não apenas o impacto das defasagens climáticas, mas
também essas características permanentes, produzindo estimativas viesadas
\citep{wooldridge2010}.

Para contornar esse problema, introduzem-se no modelo os \textbf{efeitos fixos
individuais} $\alpha_i$, que representam um intercepto específico para cada unidade
$i$. Esses termos absorvem todos os fatores não observados que diferenciam uma
unidade das demais e que não variam ao longo do tempo, como localização geográfica,
estrutura produtiva, nível histórico de renda ou composição setorial da carteira de
crédito. Em termos práticos, a estimação passa a explorar apenas a variação interna
de cada unidade ao longo do tempo, comparando cada estado consigo mesmo em períodos
distintos, em vez de compará-lo diretamente com outros estados.

Também é comum incluir \textbf{efeitos fixos temporais} $\lambda_t$, isto é, um
parâmetro específico para cada período $t$. Esses efeitos capturam choques comuns
a todas as unidades em determinado momento, como alterações na taxa básica de juros,
mudanças regulatórias nacionais, crises macroeconômicas ou períodos de expansão da
atividade econômica. Assim, controla-se tanto a heterogeneidade permanente entre as
unidades quanto fatores agregados que afetam simultaneamente toda a amostra.

Com a inclusão desses componentes, o modelo de defasagens distribuídas com efeitos
fixos pode ser escrito como

\begin{equation}
    Y_{i,t} \;=\; \alpha_i \;+\; \lambda_t \;+\;
    \sum_{k=0}^{K} \beta_k X_{i,t-k} \;+\; \varepsilon_{i,t},
    \label{eq:dlm_ef}
\end{equation}

\noindent em que $Y_{i,t}$ representa a variável dependente da unidade $i$ no
período $t$; $\alpha_i$ são os efeitos fixos individuais; $\lambda_t$ são os
efeitos fixos temporais; $X_{i,t-k}$ corresponde à variável explicativa observada
com $k$ períodos de defasagem; $\beta_k$ mede o efeito marginal associado a cada
defasagem; e $\varepsilon_{i,t}$ é o termo de erro aleatório do modelo.

\subsection{Dependência entre as Variáveis Defasadas}

Uma característica estrutural dos modelos de defasagens distribuídas é que os regressores
$X_t, X_{t-1}, \ldots, X_{t-K}$ não são variáveis independentes entre si. Ao contrário,
são versões temporalmente deslocadas de uma mesma série, e portanto tendem a ser altamente
correlacionadas. Essa correlação é conhecida como \textbf{multicolinearidade} e constitui
um dos principais desafios práticos na estimação de DLMs \citep{judge1985}.

A principal consequência da multicolinearidade é o aumento da variância dos
estimadores individuais $\hat{\beta}_k$. Em outras palavras, mesmo que os
coeficientes permaneçam não viesados sob as hipóteses usuais do modelo, seus erros
padrão tendem a se elevar, reduzindo a precisão das estimativas. Com isso, torna-se
mais difícil identificar separadamente o efeito de cada defasagem, e os testes de
significância individuais podem se tornar pouco informativos. Em situações mais
intensas, pequenas alterações na amostra ou na especificação do modelo podem gerar
mudanças relevantes nos coeficientes estimados, ainda que a associação acumulada no
horizonte $K$ $\hat{B}_K$ permaneça relativamente estável.

Para lidar com esse problema, a literatura propõe duas abordagens principais. A primeira é o
\textbf{estimador de Almon} \citep{almon1965}, que impõe a restrição de que os coeficientes
$\beta_k$ seguem um polinômio de grau $q < K$, especificados como 

\begin{equation}
    \beta_k \;=\; \sum_{j=0}^{q} a_j\, k^j, \qquad k = 0, 1, \ldots, K,
    \label{eq:almon}
\end{equation}

\noindent substituindo os $K+1$ parâmetros livres por apenas $q+1$ parâmetros $a_j$ a estimar.
Essa restrição reduz a dimensão do perfil de defasagens e pode melhorar a precisão, sem garantir menor multicolinearidade,
ao custo de impor uma forma funcional específica ao perfil. A segunda abordagem é o
\textbf{modelo de Koyck} \citep{koyck1954}, que assume decaimento geométrico dos coeficientes, $\beta_k = \beta_0 \lambda^k$, com $0 < \lambda < 1$, transformando o DLM em um
modelo autorregressivo de defasagem distribuída (ARDL) com apenas dois parâmetros.
Ambas as abordagens trocam flexibilidade por parcimônia e redução da variância dos
estimadores, sendo a escolha entre elas guiada pelo conhecimento substantivo sobre o
padrão esperado do perfil de defasagens.

A derivação formal do \textbf{estimador de Almon} torna mais preciso o
mecanismo de redução de parâmetros descrito acima. Para ilustrá-la,
considere o caso com $K = 4$ defasagens e polinômio de grau $q = 3$.
O modelo de partida é
 
\begin{equation}
    Y_t \;=\; \alpha
             + \beta_0 X_t
             + \beta_1 X_{t-1}
             + \beta_2 X_{t-2}
             + \beta_3 X_{t-3}
             + \beta_4 X_{t-4}
             + \varepsilon_t.
    \label{eq:dlm_k4_almon}
\end{equation}
 
A restrição polinomial de grau $q = 3$ impõe que cada $\beta_k$ seja
expresso como função da defasagem $k$, ou seja, 
 
\begin{equation}
    \beta_k \;=\; a_0 + a_1 k + a_2 k^2 + a_3 k^3,
    \qquad k = 0,1,2,3,4,
    \label{eq:almon_poly}
\end{equation}

\noindent de modo que os cinco coeficientes $\{\beta_k\}$ passam a ser
determinados por apenas quatro parâmetros $\{a_0, a_1, a_2, a_3\}$.
Substituindo $k = 0, 1, 2, 3, 4$ em~\eqref{eq:almon_poly}
obtém-se explicitamente:
 
\begin{align}
    \beta_0 &= a_0, \label{eq:b0} \\
    \beta_1 &= a_0 + a_1 + a_2 + a_3, \label{eq:b1} \\
    \beta_2 &= a_0 + 2a_1 + 4a_2 + 8a_3, \label{eq:b2} \\
    \beta_3 &= a_0 + 3a_1 + 9a_2 + 27a_3, \label{eq:b3} \\
    \beta_4 &= a_0 + 4a_1 + 16a_2 + 64a_3. \label{eq:b4}
\end{align}
 
Substituindo~\eqref{eq:b0}--\eqref{eq:b4} em~\eqref{eq:dlm_k4_almon}
e reagrupando os termos segundo cada parâmetro $a_j$, temos
 
\begin{align}
    Y_t \;=\; \alpha
    &+ a_0\!\left(X_t + X_{t-1} + X_{t-2} + X_{t-3} + X_{t-4}\right)
    \notag \\
    &+ a_1\!\left(X_{t-1} + 2X_{t-2} + 3X_{t-3} + 4X_{t-4}\right)
    \notag \\
    &+ a_2\!\left(X_{t-1} + 4X_{t-2} + 9X_{t-3} + 16X_{t-4}\right)
    \notag \\
    &+ a_3\!\left(X_{t-1} + 8X_{t-2} + 27X_{t-3} + 64X_{t-4}\right)
    + \varepsilon_t.
    \label{eq:almon_expanded}
\end{align}
 
O modelo em~\eqref{eq:almon_expanded} tem a estrutura de uma regressão
linear ordinária. Definindo as \textbf{variáveis auxiliares de Almon} por
 
\begin{equation}
    Z_t^{(j)} \;=\; \sum_{k=0}^{K} b_j(k)\, X_{t-k},
    \qquad j = 0, 1, \ldots, q,
    \label{eq:almon_aux_geral}
\end{equation}

\noindent
com $b_0(k)=1$ para todo $k$ e $b_j(k)=k^j$ para $j\geq1$, evitando a ambiguidade de $0^0$. O modelo pode ser reescrito de forma compacta como
 
\begin{equation}
    Y_t \;=\; \alpha
             + a_0 Z_t^{(0)}
             + a_1 Z_t^{(1)}
             + a_2 Z_t^{(2)}
             + a_3 Z_t^{(3)}
             + \varepsilon_t.
    \label{eq:almon_z}
\end{equation}
 
Na forma matricial, tem-se $\mathbf{Y} = \mathbf{Z}\,\boldsymbol{a}
+ \boldsymbol{\varepsilon}$, em que $\mathbf{Z}$ é a matriz $n\times(q+2)$
cujas colunas correspondem ao intercepto e às variáveis
$Z_t^{(0)}, \ldots, Z_t^{(q)}$, e
$\boldsymbol{a} = (\alpha, a_0, a_1, \ldots, a_q)^{\!\top}$.
O estimador de Mínimos Quadrados Ordinários fornece
 
\begin{equation}
    \hat{\boldsymbol{a}}
    \;=\;
    \bigl(\mathbf{Z}^{\!\top}\mathbf{Z}\bigr)^{-1}
    \mathbf{Z}^{\!\top}\mathbf{Y}.
    \label{eq:almon_ols}
\end{equation}
 
Obtidos os $\hat{a}_j$, os coeficientes originais do perfil de
defasagens são recuperados pela relação polinomial, dada por
 
\begin{equation}
    \hat{\beta}_k
    \;=\;
    \sum_{j=0}^{q} \hat{a}_j\, k^j,
    \qquad k = 0, 1, \ldots, K.
    \label{eq:almon_recovery}
\end{equation}
 
A transformação de Almon consiste, portanto, em projetar o vetor
$(\beta_0, \ldots, \beta_K)$ sobre o subespaço de dimensão $q + 1$
gerado pelas potências $1, k, k^2, \ldots, k^q$ avaliadas nos pontos
$k = 0, 1, \ldots, K$, a chamada matriz de Vandermonde. A redução do número de parâmetros pode melhorar a precisão quando a restrição é adequada. Entretanto, as potências brutas e as variáveis auxiliares não garantem menor correlação nem melhor condicionamento numérico; essas propriedades precisam ser diagnosticadas. O custo dessa parcimônia é a imposição
de uma forma funcional específica ao perfil de defasagens e o estimador
é não viesado somente se o grau $q$ for adequadamente especificado.
Caso $q$ seja subespecificado, introduz-se viés de restrição; caso
seja excessivamente elevado, os ganhos de variância diminuem. A
escolha de $q$ deve portanto combinar conhecimento substantivo sobre
o padrão de transmissão em análise com critérios estatísticos, como
testes de restrição ou critérios de informação AIC e BIC
\citep{judge1985, wooldridge2010}. Neste trabalho, Almon é apresentado como alternativa teórica; a aplicação utiliza o DLM finito irrestrito, sem restrição polinomial.


\subsection{Especificação e limites da aplicação empírica}
Na aplicação às 27 UFs, estima-se $y_{it}=\alpha_i+\lambda_t+\sum_{k=0}^{K}\beta_kx_{i,t-k}+\varepsilon_{it}$ com as transformações acima, efeitos fixos de UF e de mês-ano completos, sem defasagens de $y$. Os efeitos de UF representam diferenças médias nas mudanças logarítmicas, e os efeitos temporais capturam choques comuns nessa escala; não se confundem com interceptos da taxa em níveis. A identificação utiliza a variação remanescente após retirar ambos os efeitos.

Para cada exposição e recorte, escolhe-se o menor AIC entre $K=1,\ldots,24$, incluindo o contemporâneo e mantendo a amostra comum admissível a $K=24$, inclusive na estimação final. A escolha pode variar entre desastres e períodos, mas não comprova a duração física dos impactos. Os testes, intervalos pontuais e ajustes BH usuais tratam o horizonte escolhido como fixo: não incorporam a incerteza da seleção e não constituem inferência seletiva formal. BH trata a multiplicidade nas famílias declaradas, sem corrigir a seleção de $K$ ou assegurar controle global entre todas as famílias. As rejeições permanecem exploratórias.

A regra de suporte utiliza contagens contemporâneas na janela original e nas datas das equações comuns. Isso é distinto do suporte das regressoras defasadas: Onda de Frio no pré foi bloqueada pela regra de suporte (quatro UFs contemporaneamente expostas), embora as regressoras até $K=24$ alcancem seis UFs e tenham posto completo. Não se demonstrou impossibilidade de estimação em todos os horizontes, nem se assegura inferência confiável pelo posto completo. Preserva-se a regra original, sem relaxamento orientado pelos resultados.

A inferência compara cluster por UF e Driscoll--Kraay com banda temporal fixa de seis meses, independente de $K$. Ter 27 UFs não implica 27 clusters igualmente informativos. Covariâncias conjuntas mal condicionadas, sobretudo em Calor/Baixa Umidade e Chuvas Intensas, exigem cautela com testes Wald de muitas restrições; valores muito pequenos de $p$ não constituem evidência conclusiva. Essas correções não resolvem endogeneidade, seleção do horizonte ou raiz unitária.

Mantém-se a especificação transformada nos casos de diagnóstico inconclusivo em que KPSS não rejeita sua hipótese nula, registrando a limitação. A não rejeição de KPSS não prova estacionariedade, e os casos desfavoráveis também permanecem documentados (RO/SC no pré e RO no total para $\Delta\log(I)$). Os diagnósticos estaduais não certificam o painel nem a subamostra truncada; a análise permanece exploratória.

\subsection{Passo a Passo para Aplicação}

A aplicação dos modelos de defasagens distribuídas pode ser organizada nas etapas
descritas a seguir:

\textbf{1. Definição do horizonte de defasagens $K$.} O primeiro passo consiste em
determinar o número máximo de defasagens $K$ a ser incluído no modelo. Essa escolha
deve conciliar fundamentos teóricos e critérios empíricos. Do ponto de vista
econômico, $K$ deve ser suficientemente amplo para captar o intervalo de tempo
plausível entre a ocorrência do choque climático e seus possíveis efeitos sobre a
inadimplência. Do ponto de vista estatístico, critérios de informação, como AIC e
BIC, podem auxiliar na seleção de uma especificação mais parcimoniosa
\citep{wooldridge2010};

\textbf{2. Verificação da estacionaridade.} Antes da estimação, recomenda-se avaliar
se as séries utilizadas apresentam comportamento estacionário, isto é, se suas
propriedades estatísticas permanecem relativamente estáveis ao longo do tempo. Para
isso, podem ser empregados testes de raiz unitária, como o teste ADF
\citep{dickey1979} e o teste de Phillips-Perron \citep{phillips1988}. A presença de
não estacionaridade pode levar a regressões espúrias, comprometendo a interpretação
dos coeficientes e a validade das inferências \citep{greene2012};

\textbf{3. Estimação do modelo.} Definida a estrutura de defasagens e verificadas as propriedades das séries, estima-se o modelo~\eqref{eq:dlm_ef} por Mínimos Quadrados Ordinários (MQO) com efeitos fixos individuais e temporais. Os efeitos fixos individuais eliminam a heterogeneidade permanente entre unidades, ao passo que os efeitos fixos temporais removem choques agregados comuns a todos os estados em cada período. Com isso, os coeficientes $\beta_k$ são identificados exclusivamente a partir das variações temporais observadas dentro de cada unidade analisada \citep{wooldridge2010};

\textbf{4. Inferência com erros padrão robustos.} Após a estimação, é necessário
ajustar a inferência estatística para possíveis violações das suposições clássicas do
modelo. Em painéis, os resíduos podem apresentar autocorrelação ao longo do tempo
dentro da mesma unidade e heterocedasticidade entre unidades. Por esse motivo,
utilizam-se erros padrão robustos e clusterizados no nível da unidade $i$
\citep{wooldridge2010}. Essa correção torna os testes de hipótese e intervalos de
confiança mais confiáveis nessas situações;

\textbf{5. Análise do perfil de resposta defasada.} Por fim, interpretam-se os
coeficientes estimados $\hat{\beta}_0, \hat{\beta}_1, \ldots, \hat{\beta}_K$ e seus
respectivos intervalos de confiança, geralmente por meio de representação gráfica.
Esse conjunto de coeficientes forma o perfil de resposta defasada, que permite
identificar: (i) em qual período o efeito é mais intenso; (ii) se o impacto ocorre
de forma imediata ou com atraso; (iii) se os efeitos são transitórios ou
persistentes; e (iv) a associação acumulada no horizonte $K$ $\hat{B}_K$, obtido
pela soma apresentada em~\eqref{eq:efeito_lp}.

\subsection{Suposições e Limitações}

A validade dos modelos de defasagens distribuídas depende do atendimento de algumas
suposições teóricas e econométricas. Além disso, como qualquer ferramenta empírica,
esses modelos apresentam limitações que devem ser consideradas na interpretação dos
resultados.

A primeira suposição é a de \textbf{linearidade}. Assume-se que o efeito de cada
defasagem de $X$ sobre $Y$ ocorre de forma linear e aditiva, isto é, que variações
em $X$ produzem impactos proporcionais sobre $Y$, mantidos constantes os demais
fatores. Dessa forma, a especificação básica pode não captar adequadamente relações
não lineares, como efeitos de limiar, assimetrias ou mudanças de intensidade após
determinados níveis da variável explicativa.

A segunda refere-se à \textbf{exogeneidade estrita} da variável explicativa. Para
que os estimadores obtidos por MQO sejam não viesados, requer-se que $X_{t-k}$ não
esteja correlacionado com o termo de erro $\varepsilon_t$ para todo $k$
\citep{wooldridge2010}. No contexto deste trabalho, essa hipótese é plausível, pois
os desastres ambientais decorrem majoritariamente de fatores climáticos e
geofísicos externos ao sistema financeiro. Ainda assim, a existência de variáveis
omitidas correlacionadas simultaneamente com os eventos climáticos e com a
inadimplência pode comprometer essa condição.

A terceira limitação está associada à \textbf{multicolinearidade} entre as
defasagens da variável explicativa, tema discutido na seção anterior. Como valores
sucessivos de uma mesma série costumam ser fortemente correlacionados, torna-se mais
difícil separar com precisão o efeito individual de cada $\beta_k$. Em consequência,
os erros padrão podem se elevar e os coeficientes estimados tornam-se mais sensíveis
a pequenas alterações na especificação \citep{judge1985}.

A quarta questão envolve a \textbf{escolha do número de defasagens} $K$. Se o modelo
inclui poucas defasagens, efeitos relevantes podem ser omitidos, gerando viés por
variável omitida. Por outro lado, a inclusão excessiva de defasagens reduz os graus
de liberdade, diminui a precisão das estimativas e pode intensificar a
multicolinearidade. Por isso, a definição de $K$ deve combinar critérios
estatísticos com conhecimento substantivo sobre o mecanismo de transmissão em
análise.

Por fim, assim como ocorre no teste de causalidade de Granger, modelos de
defasagens distribuídas identificam principalmente relações de
\textbf{associação temporal}, e não necessariamente relações causais em sentido
estrito. A interpretação dos coeficientes $\beta_k$ como efeitos causais depende da
plausibilidade das hipóteses de identificação adotadas, especialmente da
exogeneidade da variável explicativa em relação ao termo de erro.


\section{SARIMAX}

A terceira etapa metodológica consiste na estimação de modelos SARIMAX (\textit{Seasonal Autoregressive Integrated Moving Average with Exogenous Variables}), que estendem a classe de modelos SARIMA ao incorporar variáveis exógenas à estrutura autorregressiva e de médias móveis da série temporal \citep{box2015}. Essa abordagem é particularmente adequada ao problema em análise, uma vez que a inadimplência apresenta padrões de sazonalidade e autocorrelação que devem ser adequadamente modelados antes da avaliação do efeito incremental de fatores climáticos sobre sua dinâmica.

O objetivo central dessa etapa é verificar se as variáveis climáticas agregam poder preditivo à modelagem da inadimplência além daquele já capturado pelo próprio histórico da série. Para isso, adota-se uma estratégia de comparação entre dois modelos estimados sobre o mesmo período de treinamento. O Modelo 1, utilizado como referência (\textit{baseline}), consiste em um SARIMA univariado, especificado exclusivamente a partir da dinâmica passada da inadimplência. O Modelo 2, por sua vez, corresponde ao SARIMAX, incorporando variáveis climáticas como regressores exógenos, sendo representado por

\begin{equation}
    \Phi(B)\tilde{\Phi}(B^s)\Delta^d\Delta_s^D \, y_t = \Theta(B)\tilde{\Theta}(B^s)\varepsilon_t + \sum_{j} \delta_j x_{j,t} ,
\end{equation}

\noindent em que $B$ é o operador de defasagem; $\Phi(B)$ e $\Theta(B)$ representam os polinômios autorregressivo e de médias móveis não sazonais; $\tilde{\Phi}(B^s)$ e $\tilde{\Theta}(B^s)$ capturam os componentes sazonais de periodicidade $s$; $\Delta^d$ e $\Delta_s^D$ correspondem aos operadores de diferenciação regular e sazonal; e $x_{j,t}$ denota as variáveis climáticas exógenas associadas aos coeficientes $\delta_j$. Os operadores, em especial, serão melhores explicados a seguir. 

\subsection{Processo ARMA e suas Extensões}
\label{sec:arma}
 
A construção do modelo SARIMAX parte da classe de processos
\textbf{ARMA}(\(p,\,q\)) (\textit{Autoregressive Moving Average}),
formulada por \citet{box2015} como a combinação de um componente
autorregressivo de ordem \(p\) com um componente de médias móveis de ordem
\(q\). Dado que \(y_t\) é uma série temporal estacionária com média zero, o
processo ARMA(\(p,\,q\)) é definido por
 
\begin{equation}
  y_t \;=\; \sum_{i=1}^{p} \phi_i\, y_{t-i}
           \;+\; \varepsilon_t
           \;+\; \sum_{j=1}^{q} \theta_j\, \varepsilon_{t-j},
  \label{eq:arma}
\end{equation}
 
\noindent
em que \(\phi_i\) são os coeficientes autorregressivos, \(\theta_j\) são os
coeficientes de médias móveis, e \(\varepsilon_t \sim \mathcal{N}(0,\,\sigma^2)\)
é um processo de ruído branco, não correlacionado ao longo do tempo.
 
Utilizando o operador de defasagem \(B\), em que \(B^k\,y_t = y_{t-k}\),
os polinômios \(\Phi(B) = 1 - \phi_1 B - \cdots - \phi_p B^p\) e
\(\Theta(B) = 1 + \theta_1 B + \cdots + \theta_q B^q\) permitem reescrever
o modelo de forma compacta como
 
\begin{equation}
  \Phi(B)\,y_t \;=\; \Theta(B)\,\varepsilon_t.
  \label{eq:arma_compact}
\end{equation}
 
Quando a série \(y_t\) é não estacionária, situação frequente em dados
financeiros e macroeconômicos, aplica-se diferenciação de ordem \(d\)
para induzi-la à estacionaridade. Denominando \(\Delta^d = (1-B)^d\) o
operador de diferenciação regular, obtém-se a classe
\textbf{ARIMA}(\(p,\,d,\,q\)), definida como 
 
\begin{equation}
  \Phi(B)\,\Delta^d\,y_t \;=\; \Theta(B)\,\varepsilon_t.
  \label{eq:arima}
\end{equation}
 
% ----------------------------------------------------------
\subsection{Extensão Sazonal: o Modelo SARIMA}
\label{sec:sarima}
 
Séries de frequência mensal, como a taxa de inadimplência, frequentemente
exibem padrões que se repetem em intervalos regulares de \(s\) períodos,
fenômeno denominado \textbf{sazonalidade}. Para acomodar esse comportamento,
\citet{box2015} propõem a incorporação de operadores sazonais análogos aos
não sazonais. Definindo \(\tilde{\Phi}(B^s) = 1 - \tilde{\phi}_1 B^s - \cdots - \tilde{\phi}_P B^{Ps}\)
e \(\tilde{\Theta}(B^s) = 1 + \tilde{\theta}_1 B^s + \cdots + \tilde{\theta}_Q B^{Qs}\)
como os polinômios autorregressivo e de médias móveis sazonais de ordens
\(P\) e \(Q\), respectivamente, e \(\Delta_s^D = (1-B^s)^D\) o operador de
diferenciação sazonal de ordem \(D\), o modelo
\textbf{SARIMA}(\(p,\,d,\,q\))\(\times\)(\(P,\,D,\,Q\))\(_s\) é definido
por
 
\begin{equation}
  \Phi(B)\,\tilde{\Phi}(B^s)\,\Delta^d\,\Delta_s^D\,y_t
  \;=\;
  \Theta(B)\,\tilde{\Theta}(B^s)\,\varepsilon_t,
  \label{eq:sarima}
\end{equation}
 
\noindent
em que o lado esquerdo captura a dinâmica autorregressiva regular e sazonal da
série diferenciada, e o lado direito representa a estrutura de médias móveis
regular e sazonal do processo de erros. Para séries mensais, adota-se
\(s = 12\).
 
% ----------------------------------------------------------
\subsection{Incorporação de Variáveis Exógenas: o Modelo SARIMAX}
\label{sec:sarimax_form}
 
A extensão do SARIMA para o \textbf{SARIMAX} é obtida pela adição de um
vetor de \(k\) variáveis explicativas exógenas \(\mathbf{x}_t = (x_{1,t},
\ldots, x_{k,t})^\top\) com coeficientes \(\boldsymbol{\delta} = (\delta_1,
\ldots, \delta_k)^\top\), de modo que o modelo completo torna-se
 
\begin{equation}
  \Phi(B)\,\tilde{\Phi}(B^s)\,\Delta^d\,\Delta_s^D\,y_t
  \;=\;
  \Theta(B)\,\tilde{\Theta}(B^s)\,\varepsilon_t
  \;+\; \sum_{j=1}^{k} \delta_j\,x_{j,t},
  \label{eq:sarimax}
\end{equation}
 
\noindent
em que \(B\) é o operador de defasagem; \(\Phi(B)\) e \(\Theta(B)\) representam
os polinômios autorregressivo e de médias móveis não sazonais; \(\tilde{\Phi}(B^s)\)
e \(\tilde{\Theta}(B^s)\) capturam os componentes sazonais de periodicidade
\(s\); \(\Delta^d\) e \(\Delta_s^D\) correspondem aos operadores de
diferenciação regular e sazonal; e \(x_{j,t}\) denota as variáveis climáticas
exógenas associadas aos coeficientes \(\delta_j\).
 
No contexto desta pesquisa, as variáveis exógenas \(x_{j,t}\) correspondem
a indicadores de desastres naturais registrados no período \(t\). Os
coeficientes \(\delta_j\) medem o efeito marginal de cada variável climática
sobre a taxa de inadimplência, mantidos constantes os padrões autorregressivos
e de médias móveis da série. A interpretação causal desses coeficientes,
contudo, está condicionada à exogeneidade das variáveis climáticas, suposição
plausível neste contexto dado que desastres naturais derivam de processos
geofísicos e meteorológicos externos ao sistema financeiro.
 
% ----------------------------------------------------------
\subsection{Identificação das Ordens do Modelo}
\label{sec:sarimax_id}
 
A identificação das ordens \((p,\,d,\,q)\) e \((P,\,D,\,Q)_s\) segue a
metodologia Box-Jenkins \citep{box2015}, que compreende três etapas
sequenciais: identificação, estimação e verificação diagnóstica.
 
Na etapa de identificação, as funções de autocorrelação (FAC) e de
autocorrelação parcial (FACP) da série, após as diferenciações necessárias
para eliminar a não estacionaridade regular e sazonal, são inspecionadas
para sugerir candidatos às ordens do modelo. Padrões de decaimento
exponencial na FAC com corte abrupto na FACP na defasagem \(p\) indicam
um processo AR(\(p\)); o padrão inverso, com corte na FAC na defasagem
\(q\) e decaimento na FACP, sugere um MA(\(q\)). Padrões mistos são
associados a processos ARMA.
 
Para as ordens sazonais, o mesmo raciocínio é aplicado aos \textit{lags}
múltiplos de \(s = 12\): a FAC e a FACP avaliadas nas defasagens \(s, 2s,
3s, \ldots\) orientam a escolha de \(P\) e \(Q\).
 
Ao longo dessa etapa, a seleção entre modelos concorrentes é guiada pelos
critérios de informação de Akaike \citep[AIC;][]{akaike1974} e de Schwarz
\citep[BIC;][]{schwarz1978}, definidos respectivamente por
 
\begin{equation}
  \mathrm{AIC} \;=\; -2\,\ell(\hat{\boldsymbol{\psi}}) + 2\,m,
  \label{eq:aic}
\end{equation}
 
\begin{equation}
  \mathrm{BIC} \;=\; -2\,\ell(\hat{\boldsymbol{\psi}}) + m\,\ln(T),
  \label{eq:bic}
\end{equation}
 
\noindent
em que \(\ell(\hat{\boldsymbol{\psi}})\) é o valor maximizado da função de
log-verossimilhança do modelo estimado, \(m\) é o número de parâmetros
livres e \(T\) é o tamanho da amostra. Ambos os critérios penalizam a
complexidade do modelo, favorecendo especificações mais parcimoniosas. O BIC
impõe penalidade mais severa e tende a selecionar modelos com menor número
de parâmetros quando a amostra é grande.
 
Na prática, adota-se um procedimento de busca em grade (\textit{grid
search}) sobre o espaço de ordens candidatas, estimando-se todos os modelos
factíveis e selecionando aquele com menor valor de AIC (ou BIC). Essa
estratégia reduz o risco de má especificação decorrente de inspeção visual
subjetiva das funções de autocorrelação.
 
% ----------------------------------------------------------
\subsection{Estimação}
\label{sec:sarimax_est}
 
Os parâmetros do modelo SARIMAX — os coeficientes \(\phi_i\), \(\tilde\phi_P\),
\(\theta_j\), \(\tilde\theta_Q\) e \(\delta_j\), bem como a variância dos
erros \(\sigma^2\) — são estimados pelo método de \textbf{Máxima
Verossimilhança} (MV). Sob a hipótese de normalidade dos erros
\(\varepsilon_t \sim \mathcal{N}(0,\,\sigma^2)\), a função de
log-verossimilhança do modelo tem a forma
 
\begin{equation}
  \ell(\boldsymbol{\psi}) \;=\;
    -\,\frac{T}{2}\ln(2\pi\sigma^2)
    \;-\; \frac{1}{2\sigma^2}\sum_{t=1}^{T}\varepsilon_t^2(\boldsymbol{\psi}),
  \label{eq:loglik}
\end{equation}
 
\noindent
em que \(\boldsymbol{\psi}\) denota o vetor de todos os parâmetros do modelo e
\(\varepsilon_t(\boldsymbol{\psi})\) são os resíduos condicionais obtidos
pela recursão do filtro ARMA. Na prática, a maximização de~\eqref{eq:loglik}
é realizada numericamente por meio de algoritmos de otimização não linear,
como o método BFGS, implementados em pacotes computacionais de uso
consolidado \citep{shumway2017}.
 

Após a estimação, a adequação do modelo é verificada por meio da análise
dos resíduos \(\hat\varepsilon_t\). Para que as inferências sobre os
coeficientes estimados sejam válidas, os resíduos devem se aproximar de
ruído branco, isto é, exibir média zero, variância constante e ausência de
autocorrelação serial.
 
A hipótese de ausência de autocorrelação é testada formalmente pelo
\textbf{teste de Ljung-Box} \citep{ljung1978}. A estatística de teste é
dada por
 
\begin{equation}
  Q(m) \;=\; T(T+2)\sum_{k=1}^{m}\frac{\hat{\rho}_k^2}{T - k},
  \label{eq:ljungbox}
\end{equation}
 
\noindent
em que \(\hat{\rho}_k\) é a autocorrelação amostral dos resíduos na defasagem
\(k\) e \(m\) é o número de defasagens testadas. Sob a hipótese nula de
ausência de autocorrelação, \(Q(m)\) segue assintoticamente uma distribuição
qui-quadrado com \(m - (p + q + P + Q)\) graus de liberdade. A rejeição de
\(H_0\) indica que o modelo deixou autocorrelação não captada nos resíduos,
sinalizando a necessidade de reespecificação.
 
Adicionalmente, inspecionam-se a FAC e a FACP dos resíduos e os gráficos de
diagnóstico padrão — histograma dos resíduos para avaliar a aderência à
distribuição normal, e gráfico quantil-quantil (Q-Q plot) para detectar
desvios nas caudas da distribuição.
 
% ----------------------------------------------------------
\subsection{Avaliação Preditiva e Comparação entre Modelos}
\label{sec:sarimax_eval}
 
A avaliação comparativa entre os modelos é conduzida fora da amostra, por
meio de um conjunto de validação não utilizado na etapa de estimação. A
capacidade preditiva é mensurada a partir de métricas usuais de erro de
previsão: a raiz do erro quadrático médio (RMSE), que penaliza desvios de
maior magnitude; o erro absoluto médio (MAE), mais robusto à presença de
valores extremos; e o erro percentual absoluto médio (MAPE), que fornece
uma medida relativa do erro. Adicionalmente, os critérios de informação AIC
e BIC são empregados na seleção das ordens dos modelos, favorecendo
especificações parcimoniosas.
 
Caso o modelo SARIMAX apresente desempenho superior ao modelo SARIMA nas
métricas de previsão, isso constitui evidência de que as variáveis climáticas
possuem poder preditivo incremental sobre a inadimplência. Em outras
palavras, tal resultado sugere que a ocorrência de desastres ambientais
contribui para explicar a dinâmica da inadimplência além do que pode ser
inferido exclusivamente a partir de seu comportamento passado.